% !TEX root = ../../../tree_notes.tex
% =====================================================================
%  14.4  Games on trees: Grundy numbers and win/lose flags   sec:c27-games
%  Exports: def:c27-grundy  thm:c27-grundy-xor  thm:c27-hackenbush
%           lst:c27-grundy  fig:c27-gamefig  ex:c27-grundy-run
%           prob:c27-game-solved  lst:c27-455b-soln
%  The last budgeted DP of the chapter has an ADVERSARY instead of a knapsack:
%  two players move on a tree/trie and we ask who is forced to lose.  Two
%  tools: the Sprague-Grundy number g(v) = mex of the option values, whose
%  disjoint sums combine by XOR (thm:c27-grundy-xor); and, for a cutting game,
%  the Green-Hackenbush identity g(v) = xor over children of (g(c)+1)
%  (thm:c27-hackenbush).  Hand-verified on the running tree T: g(1) = 7.
%  Anchor CF 455B (A Lot of Games): a two-flag win/lose DP on a trie, samples
%  First / First / Second re-derived from the flags; k-game parity logic.
%  Practice CF 1153D (Serval and Rooted Tree): minimax tree DP, samples
%  1, 4, 4, 5 re-derived from leaf=1 / max->min / min->sum / k+1-dp[root].
% =====================================================================
\section{Games on Trees: Grundy Numbers and Win/Lose Flags}\label{sec:c27-games}

\begin{previouslybox}
\begin{enumerate}[nosep, leftmargin=1.6em]
  \item In \secref{c27-knapsack} the fold combined children with
        $\max\!\sum$ (a shared budget). If instead each child were an
        \emph{independent sub-game} and a move chose one child to play in, what
        operation combines the children now? (not a sum -- an
        $\bigoplus$ of nim-values, \thmref{c27-grundy-xor})
  \item Chapter~13's two-state MIS (\secref{c10-twostate}) tagged every node
        \emph{taken/not-taken}. A game DP tags a node by who wins from it.
        What are the two cheapest labels, before any Grundy number? (a
        $\mathit{win}$ flag and a $\mathit{lose}$ flag)
  \item \secref{c27-greedy} proved a greedy optimal by arguing the local choice
        was \emph{forced}. A forced-loss position is the game analogue: if every
        move loses, what is the node's label? (a losing node -- $\mathit{win}=\mathit{false}$)
  \item The empty game (no moves for the player to move) -- is that a win or a
        loss for the player whose turn it is? (a loss: he cannot move; this is
        the leaf base case $g = 0$)
\end{enumerate}
\end{previouslybox}

\begin{topicintro}
Every section of this chapter spent a budget or a profile \emph{on your side}.
The last tree DP has an \emph{opponent}: two players move alternately along the
tree (or a trie of legal positions) and you must decide who is \emph{forced} to
run out of moves. The generic engine is the \term{Sprague-Grundy number}: to
each position $v$ you attach $g(v) = \operatorname{mex}$ of the $g$-values of its options, and
the whole game is decided by whether $g$ at the root is zero. Two facts make
this tractable on trees. First, a position that splits into independent branches
(a move plays in exactly one branch) has value the \emph{XOR} of the branch
values -- \thmref{c27-grundy-xor}, the reason trees and nim fit together so well.
Second, for the canonical \emph{edge-cutting} game on a rooted tree (cut an edge,
the subtree above it falls), the values obey the closed form
$g(v) = \bigoplus_{c} (g(c) + 1)$ -- \thmref{c27-hackenbush} -- computable in one
post-order DFS, exactly the shape of \secref{c27-knapsack}'s fold with $\oplus$
in place of $\max$.

We then solve the section's anchor, CF~455B (\probref{c27-game-solved}), which is
\emph{not} a Grundy problem but its minimal cousin: a \emph{two-flag} win/lose DP
on a trie, and the simplest game-on-a-tree that shows the ``forced to move,
forced to lose'' paradigm; the $k$-round rule is a tiny parity argument on top.

\textbf{Prerequisites.} \secref{c27-knapsack}'s post-order child fold (reused
verbatim, $\oplus$ replacing $\max$); Ch.~13's two-state labelling
\secref{c10-twostate}; the notion of $\operatorname{mex}$ (minimum excluded value), defined
below; Ch.~4's rooting. Full combinatorial-game theory lives elsewhere -- this
section is self-contained for the tree DP you will actually write.

\textbf{After this section you can}
\begin{itemize}[nosep]
  \item compute $\operatorname{mex}$ option sets and $g(v)$ for a pruning game by one DFS; \hfill\lv{2}--\lv{3}
  \item justify combining branches with XOR rather than addition; \hfill\lv{4}
  \item cold-write CF~455B's win/lose trie DP and its $k$-game verdict. \hfill\lv{4}--\lv{5}
\end{itemize}
\end{topicintro}

\subsection{Intuition}

A one-player DP asks ``what is the best I can do?'' A two-player, perfect-information,
no-draw game asks ``can I \emph{force} a win?'' Both collapse to a post-order fold, but
the quantifier changes: instead of $\max$ over children you get
$\mathit{win}(v) = \bigvee_c \neg\,\mathit{win}(c)$ -- ``I move somewhere my opponent
cannot win from.'' That is already enough for the anchor (CF~455B). Sprague-Grundy
theory goes further and measures a position by a whole \emph{number} $g(v)$, not a
Boolean, so that positions you have never seen can be \emph{added} (played as disjoint
sums) by arithmetic.

Two ideas carry everything. (1) The value of a position is $\operatorname{mex}$ of the values of its
options: the smallest non-negative integer the mover \emph{cannot} reach. A terminal
position has no options, so $g = \operatorname{mex}(\varnothing) = 0$. (2) A position that decomposes
into independent components -- a move touches exactly one of them -- is the
\emph{bitwise XOR} $\bigoplus$ of the components, not their sum. XOR is the operation
because it is its own inverse in the ``copy'' sense: two identical components cancel
($x \oplus x = 0$), exactly matching the strategy-stealing mirror play a second player
uses to neutralise a symmetric position. On a rooted tree the components are the child
branches, so a single DFS folds $g$ up with $\bigoplus$ -- the very shape of
\secref{c27-knapsack}, with one operator swapped.

\subsection{Definitions and Notation}

\begin{definition}{Grundy number of a tree position}{c27-grundy}
Let $v$ be a position whose options (positions reachable in one legal move) are
the nodes $u$ with an edge $v \to u$. The \term{Grundy number} (or \term{nimber})
is
\[
  g(v) \;=\; \operatorname{mex}\big(\{\, g(u) : u \text{ an option of } v \,\}\big),
  \qquad \operatorname{mex}(S) = \min\{\, t \ge 0 : t \notin S\,\}.
\]
For an \emph{impartial} game (the same moves available to both players) on a tree
directed away from the root, options of $v$ are exactly its children, so $g$ folds
bottom-up: a leaf has $g = \operatorname{mex}(\varnothing) = 0$. The single-game verdict is
$v$ is a first-player win iff $g(v) \ne 0$.
\end{definition}

\noindent\textbf{Reading $\operatorname{mex}$.} Collect the children's $g$-values into a set $S$;
$g(v)$ is the smallest non-negative integer \emph{missing} from $S$. E.g.\
$S = \{0,1,3\} \Rightarrow g = 2$; $S = \{1,2\} \Rightarrow g = 0$; $S = \varnothing
\Rightarrow g = 0$. Because you only ever need to know which small integers are present,
$\operatorname{mex}$ of a node with $d$ children is computed by a boolean scan of at most $d+1$ slots --
$\bigO{d}$ -- so the whole tree is $\bigO{n}$.

\begin{theorem}{Disjoint sums combine by XOR (Sprague-Grundy)}{c27-grundy-xor}
Let a position $P$ be the disjoint sum of independent components $P_1, \dots, P_m$
(a move chooses exactly one $P_i$ and moves within it). Then
$g(P) = g(P_1) \oplus g(P_2) \oplus \cdots \oplus g(P_m)$, where $\oplus$ is
bitwise exclusive-or. In particular $P$ is a losing position for the player to move
iff $g(P) = 0$.
\end{theorem}
\begin{proofbox}[\thmref{c27-grundy-xor}: the two defining properties of $\operatorname{mex}$ transfer through $\oplus$]
Set $x = \bigoplus_i g(P_i)$. To prove $g(P) = x$ we show (i) no option of $P$ has
value $x$, and (ii) every $y < x$... more precisely every option-value below is
reachable -- the standard two conditions that characterise a nimber. \emph{(i) No move
keeps the xor fixed.} Any move changes exactly one component $P_j \to P_j'$, hence
$g(P_j) \to g(P_j') \ne g(P_j)$ (a Grundy number never repeats among its own options).
Changing one operand of an XOR strictly changes the result, so no option of $P$ has
value $x$. \emph{(ii) Every value below is reachable.} To move $P$ to value $z < x$:
look at the most-significant bit where $x$ and $z$ differ (it is a $1$-bit of $x$);
some component $P_j$ has that bit set in $g(P_j)$; moving within $P_j$ to a component
of value $g(P_j) \oplus x \oplus z < g(P_j)$ is legal by the option-realisation property
of a single component. \emph{Zero test:} if $x = 0$, every move makes $x \ne 0$ (by (i))
so the mover hands the opponent a nonzero -- winnable -- position: $0$ is a loss. If
$x \ne 0$, (ii) with $z = 0$ gives a move to a losing position for the opponent: a win.
\qed
\end{proofbox}

\begin{theorem}{Green Hackenbush on a rooted tree: $g(v) = \bigoplus_c (g(c) + 1)$}{c27-hackenbush}
Play the impartial game on a tree rooted at the ground: a move cuts one edge, and the
cut subtree (everything on the far side) is removed. Then the Grundy number of the
position rooted at $v$ satisfies
\[
  g(v) \;=\; \bigoplus_{c \in \children(v)} \big(g(c) + 1\big),
\]
with $g(\ell) = 0$ at a leaf (the empty XOR). It is a single post-order fold.
\end{theorem}
\begin{proofbox}[\thmref{c27-hackenbush}: branches are a disjoint sum, and a branch is a stalk of height $g(c)+1$]
The branches hanging from $v$ are independent components: a move either cuts the edge
$(v, c)$ (deleting all of branch $c$) or is played strictly inside branch $c$, and no
move touches two branches. By \thmref{c27-grundy-xor} the position is the XOR of the
branch values, so it suffices to value one branch. A branch $=$ the single edge $(v,c)$
capped by the sub-position rooted at $c$; that edge acts like one extra nim-stalk unit
glued on top of $c$, and a stalk of $h$ edges has value $h$ (a path lets the mover cut
to any smaller length, giving $\operatorname{mex}\{0, \dots, h-1\} = h$). Gluing a unit stalk below a
position of value $g(c)$ raises its value to $g(c) + 1$ (the fusion/colon principle for
a chain of green edges). Hence branch $c$ contributes $g(c) + 1$, and summing over
children with $\oplus$ gives the claim; the leaf has no branches so $g = \bigoplus
\varnothing = 0$. \qed
\end{proofbox}

\begin{insightbox}
\thmref{c27-hackenbush} turns ``who wins the tree-cutting game on this whole rooted tree?''
into one DFS with a single XOR per node -- the cheapest possible reuse of the
\secref{c27-knapsack} template (swap $\max\!\sum$ for $\bigoplus$). Read it as the
chapter's punchline: \emph{a tree DP is any bottom-up fold; the operator depends on the
problem}. Budget $\to \max$; growing profile $\to$ the tails array; a game $\to \operatorname{mex}$ and
$\oplus$. And note the ground: for these cutting games you never need general $\operatorname{mex}$ --
the closed form already encodes it -- but you \emph{do} need $\operatorname{mex}$ the moment the move
set is not a clean subtree deletion (the anchor CF~455B is such a case: a leaf is forced
to move, so we use flags, not a stalk count).
\end{insightbox}

\subsection{The Algorithm: Grundy Fold by Post-Order DFS}

\begin{lstlisting}[caption={Generic Grundy fold (mex over children) and the Hackenbush shortcut (xor of child+1); both are one post-order pass, same skeleton as \coderef{c27-knap}}, label={lst:c27-grundy}]
// Root the tree at the ground node; children via a parent-guard DFS.
vector<vector<int>> adj;

// --- General impartial game: g[v] = mex of children's g ---
int grundy_mex(int v, int p) {
    vector<char> seen(adj[v].size() + 1, 0);   // g is <= number of children
    for (int u : adj[v]) if (u != p) seen[grundy_mex(u, v)] = 1;
    int g = 0; while (seen[g]) g++;            // mex = smallest unseen
    return g;
}

// --- Cutting game (thm:c27-hackenbush): g[v] = XOR over children of (g[c]+1) ---
int hackenbush(int v, int p) {
    int g = 0;                                 // empty XOR = 0 (leaf base)
    for (int u : adj[v]) if (u != p) g ^= (hackenbush(u, v) + 1);
    return g;
}

// Single-game verdict: the mover at the root wins iff the value is nonzero.
// First wins  <=>  hackenbush(ground, -1) != 0.
\end{lstlisting}

\begin{complexitybox}
Both folds are one post-order DFS: $\bigO{n}$ time, $\bigO{n}$ memory, the recursion
depth $\bigO{h}$ (a chain is legal, so watch the stack as in \secref{c27-knapsack}).
In \code{grundy\_mex} the \code{seen} array has length $\deg(v) + 1$ because $g(v) \le
\deg(v)$ (at most $\deg(v)$ distinct option values to skip), so the mex scan across all
nodes is $\sum_v \bigO{\deg(v)} = \bigO{n}$ -- no hidden cost. Verdict: the fold is
linear \emph{because} it is a plain bottom-up pass, exactly like Ch.~13's constant-tuple
fold; unlike \secref{c27-knapsack} there is no budget $V$ to amortise, since a node
carries a \emph{single} $g$-value. This is the ``$\bigO{n}$ even without a cap'' case.
\end{complexitybox}

\subsection{Worked Example: Grundy Overlay on the Running Tree $T$}

\begin{example}{The cutting-game value of the running tree $T$ rooted at $1$}{c27-grundy-run}
Apply \thmref{c27-hackenbush} to $T$ (rooted at $1$): each node's value is the XOR of
(child value $+ 1$). The bottom-up arithmetic gives leaves $0$ and, at the root,
$g(1) = 7 \ne 0$: the first player to cut wins, and can move to a $0$-value position.
\end{example}

\begin{figure}[H]
\centering
\begin{tikzpicture}[diagram, scale=0.82, transform shape,
    tnode/.style={circle, draw=primaryblue, line width=0.9pt, fill=primaryblue!8,
                  minimum size=7mm, font=\small\bfseries, text=primaryblue!60!black}]
  \node[tnode] (n1) at (0,0) {$1$};
  \node[tnode] (n2) at (-3,-1.6) {$2$};
  \node[tnode] (n3) at ( 3,-1.6) {$3$};
  \node[tnode] (n4) at (-5.2,-3.2) {$4$};
  \node[tnode] (n5) at (-1.6,-3.2) {$5$};
  \node[tnode] (n6) at ( 1.6,-3.2) {$6$};
  \node[tnode] (n7) at ( 5.2,-3.2) {$7$};
  \node[tnode] (n8) at (-1.6,-4.8) {$8$};
  \node[tnode] (n9) at ( 5.2,-4.8) {$9$};
  \node[tnode] (n10) at ( 5.2,-6.4) {$10$};
  \draw[tedge] (n1)--(n2) (n1)--(n3) (n2)--(n4) (n2)--(n5)
               (n3)--(n6) (n3)--(n7) (n5)--(n8) (n7)--(n9) (n9)--(n10);
  % Grundy badges (hackenbush values)
  \foreach \nd/\gg in {n1/7,n2/3,n3/2,n4/0,n5/1,n6/0,n7/2,n8/0,n9/1,n10/0} {
    \node[font=\tiny\ttfamily, text=successgreen!55!black, anchor=south west] at (\nd.south west) {\gg};}
  \node[annot, font=\scriptsize, text=successgreen!50!black, anchor=west] at (5.9,-3.2)
    {green badge $= g(v)$};
  \node[annot, font=\scriptsize, text=primaryblue!60!black, anchor=west] at (5.9,-3.9)
    {$g(1)=4 \oplus 3 = 7$};
  \node[annot, font=\scriptsize, text=difforange!75!black, anchor=west] at (5.9,-4.8)
    {nonzero $\Rightarrow$ First wins};
\end{tikzpicture}
\caption{The Hackenbush Grundy overlay of $T$, rooted at $1$. Leaves
$4,6,8,10$ have $g=0$; then $g(9)=g(10){+}1=1$, $g(5)=1$, $g(7)=g(9){+}1=2$;
$g(2)=(0{+}1)\oplus(1{+}1)=1\oplus2=3$ and $g(3)=(0{+}1)\oplus(2{+}1)=1\oplus3=2$;
finally $g(1)=(3{+}1)\oplus(2{+}1)=4\oplus3=7$. Every XOR re-checks by hand, so the
first cutter wins and can move the root to a $0$-value position.}
\label{fig:c27-gamefig}
\end{figure}

\begin{tracebox}[Bottom-up Grundy fold on $T$ ($g(v) = \bigoplus_c (g(c)+1)$)]
leaves: $g(4)=g(6)=g(8)=g(10)=0$.\\
$9$: child $10$ $\Rightarrow (0+1)=1$. \;\;$g(9)=1$.\\
$5$: child $8$ $\Rightarrow (0+1)=1$. \;\;$g(5)=1$.\\
$7$: child $9$ $\Rightarrow (1+1)=2$. \;\;$g(7)=2$.\\
$2$: children $4,5$ $\Rightarrow (0+1)\oplus(1+1)=1\oplus2=3$. \;\;$g(2)=3$.\\
$3$: children $6,7$ $\Rightarrow (0+1)\oplus(2+1)=1\oplus3=2$. \;\;$g(3)=2$.\\
$1$: children $2,3$ $\Rightarrow (3+1)\oplus(2+1)=4\oplus3=\mathbf{7}$. \checkmark\\
$7 \ne 0$ so the mover at $1$ wins: reduce the branch $2$ from $4$ to $3$ by a legal
cut inside $T_2$, which turns the top XOR to $3 \oplus 3 = 0$. \;\;Verdict: First wins.
\end{tracebox}

\begin{insightbox}[Key takeaways]
\begin{enumerate}[nosep, leftmargin=1.6em]
  \item \emph{One value per node, one XOR merge.} Games fold like \secref{c27-knapsack}
        but with $\oplus$ and no budget -- $g(v)$ is a single integer, so the pass is
        plainly $\bigO{n}$.
  \item \emph{$\operatorname{mex}$ where the move set is not a clean cut.} The Hackenbush closed form
        hides $\operatorname{mex}$; the moment a move is ``pass to any option'' rather than ``delete a
        subtree'', compute $g = \operatorname{mex}$ of the option set directly (def.~\ref{def:c27-grundy}).
  \item \emph{Zero means the mover loses.} Every verdict in this section -- the cutting
        game and the anchor below -- is ``root value $\ne 0$ $\Rightarrow$ First'', once you
        have combined branches with the \emph{right} operator.
\end{enumerate}
\end{insightbox}

\subsection{Variations and Twists}

\begin{itemize}[nosep]
  \item \textbf{Weighted edges}: if cutting edge $e$ deletes its subtree and $e$ carries
        a heap of size $w$ above it, the branch contributes $g(c) \oplus w$-like stalk
        values; \emph{bound the reachable set before claiming} $g(v) \le \height(v)$ --
        weighted moves can push $g$ past the height. \hfill\lv{4}
  \item \textbf{Not impartial}: if the two players have \emph{different} moves (partizan),
        Grundy/XOR is invalid and you need Conway games; a tree DP still exists but the
        operator is no longer $\oplus$. \hfill\lv{5}
  \item \textbf{Two-flag minimax (the anchor)}: when you only care \emph{win/lose}, not the
        number, drop to $\mathit{win}[v] = \bigvee_c \neg \mathit{win}[c]$ -- cheaper and
        exactly CF~455B. Grundy is the upgrade when positions must be \emph{added}. \hfill\lv{3}
  \item \textbf{Rerooting the game}: value the tree when the ground is each possible root is
        a rerooting fold over \thmref{c27-hackenbush} -- Ch.~15 (\chref{c20}), forward-cited. \hfill\lv{5}
\end{itemize}

\subsection{Common Pitfalls}

\begin{warningbox}[Trap 1 -- combining branches by $+$ or $\max$ instead of $\oplus$]
Disjoint sub-games compose by XOR, never by sum or max. A node with two identical
branches (each value~$2$) has value $2 \oplus 2 = 0$ (a losing position -- the second
player mirrors), but $2+2=4$ or $\max=2$ both wrongly say ``nonzero, First wins''. If
your verdict flips on a symmetric tree, you used the wrong operator.
\end{warningbox}

\begin{warningbox}[Trap 2 -- $\operatorname{mex}$ scan too short]
$g(v)$ is at most $\deg(v)$, so a $\mathit{seen}$ array of size $\deg(v)+1$ is exactly
enough; sizing it to a fixed small constant silently clips realizable values and mis-solves
high-degree nodes. Conversely do not ``bound $g$ by the height'' -- that only holds for
pure stalks, not for branching or weighted moves (see Variation~1).
\end{warningbox}

\begin{warningbox}[Trap 3 -- forgetting the leaf/base and the $+1$]
In \thmref{c27-hackenbush} each child contributes $g(c) \,\mathbf{+}\, 1$, not $g(c)$: the
$+1$ is the edge you are standing on. A leaf's value is $0$ (empty XOR). Omitting the $+1$
collapses every stalk to $0$ and reports ``Second wins'' on positions that are obvious
single-cut wins.
\end{warningbox}

\subsection{Problem Lab}

\solvedproblem{c27-game-solved}{A Lot of Games}{Codeforces 455B}
  {https://codeforces.com/problemset/problem/455/B}{CF 1700 · L4}
  {two-flag win/lose DP on a trie, minimax forced-move labels, k-round parity,
   impartial-game reasoning without Grundy numbers}

\begin{problemstatement}[A Lot of Games (statement and samples verified against the live problem)]
Given $n$ non-empty strings (total length $\le 10^5$, $n \le 10^5$). Two players build a
word from empty, alternating turns; each move appends one letter so the word stays a
prefix of \emph{at least one} string. A player who cannot move loses. They play this
\emph{single} game $k$ times ($1 \le k \le 10^9$); the loser of game $i$ moves first in
game $i+1$. The winner of the whole match is the winner of game $k$. Both play
optimally. Print \texttt{First} if the player who moves first \emph{in game~1} wins the
match, else \texttt{Second}. \textbf{Example~1}: $n=2$, $k=3$, strings \{a, b\} $\Rightarrow$
\texttt{First}. \textbf{Example~2}: $n=3$, $k=1$, \{a, b, c\} $\Rightarrow$ \texttt{First}.
\textbf{Example~3}: $n=1$, $k=2$, \{ab\} $\Rightarrow$ \texttt{Second}.
\end{problemstatement}

\begin{attemptgate}[Attempt it yourself first -- \lv{4}, give yourself 30 minutes]
Positions are the trie nodes. First get a \emph{single} game right: at a node, can the
mover force a win? Can he force a \emph{loss}? Only then fold in the $k$-game rule.
\peek{At a trie leaf the mover has no letter to add. Is that a win or a loss for him?}
\nudge{Loss -- so $\mathit{win}[\mathit{leaf}]=\mathit{false}$. For an internal node
      the mover wins iff \emph{some} child is a losing node for the opponent.}
\fullhint{Build the trie. Compute two flags per node bottom-up: $\mathit{win}[v] =
\bigvee_c \neg \mathit{win}[c]$ (can force a win) and $\mathit{lose}[v] = \bigvee_c \neg
\mathit{lose}[c]$ (can force a loss), with $\mathit{win}[\mathit{leaf}]=\mathit{false}$,
$\mathit{lose}[\mathit{leaf}]=\mathit{true}$. Read at the root: if $\neg\mathit{win}$ the
first player can never win a game $\Rightarrow$ \texttt{Second}. Else if
$\mathit{lose[\mathit{root}]}$ he can win \emph{or} lose on purpose, controlling who starts
game $k$ $\Rightarrow$ \texttt{First}. Else every game is won by its starter, so game $k$
goes to \texttt{First} iff $k$ is odd, else \texttt{Second}.}
\end{attemptgate}

\solstep{Understanding the Problem}
A single game is exactly a walk down the trie: from the current prefix node, one move =
one child. It ends at a node with no children, where the mover cannot extend the word and
loses. That is a finite impartial game, but the match is \emph{$k$} such games chained by
``loser starts next.'' So we need not just \emph{who wins one game} but \emph{whether the
first player can also choose to lose one} -- because being able to lose lets him decide who
starts the next round, and thus who starts round~$k$. Two booleans per node capture both.

\solstep{Key Observations}
\begin{observation}
$\mathit{win}[v]$ (mover can force a win) obeys $\mathit{win}[v] = \bigvee_{c}
\neg\,\mathit{win}[c]$: move to a child where the opponent cannot win. A leaf has no child
so $\mathit{win}=\mathit{false}$ -- the mover loses.
\end{observation}
\begin{observation}
$\mathit{lose}[v]$ (mover can force a loss) obeys $\mathit{lose}[v] = \bigvee_{c}
\neg\,\mathit{lose}[c]$ by the same logic on the ``lose'' objective, with a leaf giving
$\mathit{lose}=\mathit{true}$ (the mover is forced to lose). Both flags fold in one DFS.
\end{observation}
\begin{observation}
Match verdict at the root: $\neg\mathit{win}[\text{root}]$ means First loses every game
$\Rightarrow$ \texttt{Second}. $\mathit{lose}[\text{root}]$ means First may win or lose each
game at will, so he steers who starts game~$k$ and wins $\Rightarrow$ \texttt{First}.
Otherwise a win is forced for whoever starts, so the starter of game~$k$ wins; tracking
``loser starts next'' gives \texttt{First} iff $k$ is odd.
\end{observation}

\solstep{Think It Through With Me}
\begin{thinkbox}
\deadend{Compute a Grundy number $g(v)$ with $\operatorname{mex}$ and XOR the branches, as in
\thmref{c27-grundy-xor}.}
{Overkill and slightly wrong: here a move does not split the position into a disjoint
sum -- it \emph{chooses one} child and descends -- so there is nothing to XOR. The game is
a pure win/lose (a Boolean), and the flags are the minimal sufficient state. This is the
Variation~3 case: only use Grundy when positions must be added.}
\deadend{Answer is \texttt{First} iff First wins one game (ignore $k$ and the loser-starts
rule).}
{Fails Example~3: \{ab\} First \emph{can} reach a win in one game? No -- the root's only
line forces the mover's outcome; with $k=2$ the loser-starts rule flips who closes. The
$\mathit{lose}$ flag and parity exist precisely to model this. Trap~3 in a new guise.}
\aha{Two flags $\mathit{win}, \mathit{lose}$ folded bottom-up on the trie; verdict from
the pair at the root plus a parity test on $k$. The proof box below gives the case split,
and \defref{c27-grundy} places it on the Boolean-vs-Grundy ladder.}
\end{thinkbox}

\solstep{Solution}
Insert the $n$ strings into a trie (root~$0$). Post-order DFS fills
$\mathit{win}[v], \mathit{lose}[v]$. Let $W = \mathit{win}[0]$, $L = \mathit{lose}[0]$.
If $\neg W$ print \texttt{Second}; else if $L$ print \texttt{First}; else print
\texttt{First} if $k$ odd else \texttt{Second}.

\solstep{Proof of Correctness}
\begin{proofbox}[\coderef{c27-455b-soln}: the two flags and the parity rule decide the match]
\emph{Flags.} Induction on trie height. Height~$0$ (leaf): the mover cannot move, so he
loses: $W=\mathit{false}$ (no child makes $\bigvee$ true) and $L=\mathit{true}$. Step: the
mover can force a win iff some move gives the opponent a non-winning position, i.e.\
$\bigvee_c \neg W[c]$ -- the recurrence -- and identically for forcing a loss. So both flags
are correct at the root. \emph{Match.} If $\neg W$ (First cannot win game~1 as starter, and
symmetrically the starter of any game cannot win when he is the ``$\neg W$'' type), then
each game is won by its \emph{second} mover; game~$k$'s winner is Second overall:
\texttt{Second}. If $L$ also holds, First can win or lose each game he starts at will, so he
chooses the starters to make himself win game~$k$: \texttt{First}. If $W$ holds but $\neg L$
(First wins whenever he starts and cannot throw a game), the winner of every game is its
starter, and since the loser starts the next, game $k$ is started by First iff $k$ is odd:
\texttt{First} for odd $k$, \texttt{Second} for even $k$. All cases match the recurrence.
\qed
\end{proofbox}

\solstep{Implementation}
\begin{lstlisting}[caption={Codeforces 455B -- two-flag win/lose DP on the trie plus the k-round verdict (\coderef{c27-grundy} is the Grundy cousin we deliberately avoid)}, label={lst:c27-455b-soln}]
#include <bits/stdc++.h>
using namespace std;
const int MAXN = 100005;
int nxt[MAXN][26], win_[MAXN], lose_[MAXN], tot = 1; // node 0 = root

void insert(const string& s) {
    int v = 0;
    for (char ch : s) {
        int c = ch - 'a';
        if (!nxt[v][c]) nxt[v][c] = tot++;
        v = nxt[v][c];
    }
}

void dfs(int v) {
    bool any = false, canWin = false, canLose = false;
    for (int c = 0; c < 26; c++) if (nxt[v][c]) {
        any = true;
        dfs(nxt[v][c]);
        canWin  = canWin  || !win_[nxt[v][c]];
        canLose = canLose || !lose_[nxt[v][c]];
    }
    if (!any) { win_[v] = 0; lose_[v] = 1; }   // leaf: mover loses
    else { win_[v] = canWin; lose_[v] = canLose; }
}

int main() {
    ios::sync_with_stdio(false); cin.tie(0);
    int n; long long k;
    cin >> n >> k;
    for (int i = 0; i < n; i++) { string s; cin >> s; insert(s); }
    dfs(0);
    if (!win_[0])          cout << "Second\n";   // first can never win a game
    else if (lose_[0])     cout << "First\n";     // first controls who starts game k
    else cout << (k & 1 ? "First\n" : "Second\n"); // every game won by its starter
    return 0;
}
\end{lstlisting}

\begin{complexitybox}
Trie build is $\bigO{\text{total length} \cdot 1} = \bigO{26 \cdot n}$ in the array form
(or $\bigO{\text{total length}}$ with maps); the DFS visits each trie node once, $\bigO{26
\cdot \text{nodes}} = \bigO{\text{total length}}$. Memory $\bigO{26 \cdot \text{nodes}}$
($\approx 10^5 \cdot 26$ ints -- fine). $k$ is only used for a parity test, so read it in
\code{long long} ($k \le 10^9$ fits \code{int} but the flag logic does not scale with $k$).
Verdict: \ac; the recurrence is a two-Boolean game fold, the concrete ``forced-move label''
version of \defref{c27-grundy}. Recursion depth is bounded by the longest string ($\le 10^5$)
-- raise the stack or iterate if your judge caps it.
\end{complexitybox}

\begin{tracebox}[Verdict logic on Example~1: \{a, b\}, $k=3$]
Trie: root $0$ with two children $a, b$, both leaves. Leaves: $win=0,lose=1$.\\
root: $win[0] = \neg win[a] \vee \neg win[b] = \text{true}$; \;
$lose[0] = \neg lose[a] \vee \neg lose[b] = \text{false}$.\\
$win[0]$ true $\Rightarrow$ not the \texttt{Second} branch. $lose[0]$ false $\Rightarrow$ not
the \texttt{First} branch. So parity of $k=3$ (odd) $\Rightarrow$ \texttt{First}. \checkmark
\end{tracebox}

\begin{tracebox}[Examples~2 and~3, same two-flag pipeline]
Ex.~2 \{a,b,c\}, $k=1$: root has 3 leaf children, identical to Ex.~1 structurally:
$win=1, lose=0$; $k=1$ odd $\Rightarrow$ \texttt{First}. \checkmark\\
Ex.~3 \{ab\}, $k=2$: chain $0 \to a \to b(\text{leaf})$. $win[b]=0,lose[b]=1$;
$win[a]=\neg win[b]=1$, $lose[a]=\neg lose[b]=0$; $win[0]=\neg win[a]=0$. $win[0]$
false $\Rightarrow$ \texttt{Second} (regardless of $k$). \checkmark
\end{tracebox}

\solstep{Pitfalls and Variants}
\begin{warningbox}
Live traps for CF~455B: (i)~treating it as Grundy and XOR-ing children (it is a
\emph{choose-one-descend} game -- flags, not nimbers); (ii)~using $\mathit{win}$ alone and
ignoring $k$/the $\mathit{lose}$ flag, which breaks Example~3's chain; (iii)~swapping the
leaf base ($\mathit{win}=0$, $\mathit{lose}=1$) -- get these flipped and every verdict
inverts; (iv)~reading $k$ as \code{int} if you also do arithmetic on it. Contrast the
practice CF~1153D: there the two players \emph{maximise/minimise a number}, not a win, so
the recurrence is $\max/\min$ over children with a ``sacrifice count'' twist -- same
bottom-up shape, different objective.
\end{warningbox}

\subsection{Practice Problems}

\begin{stuckbox}
\begin{itemize}[nosep, leftmargin=1.4em]
  \item Is your game \emph{impartial} (same moves for both, so Grundy/$\oplus$ applies) or
        \emph{optimisation} (max/min a value, so a $\max/\min$ fold)? Name the operator
        before writing the recurrence.
  \item For a cutting game, did you add the $+1$ per child (\thmref{c27-hackenbush}) and
        use the empty-XOR leaf base $g=0$?
  \item For a win/lose flag DP, did you compute BOTH the ``can-win'' and ``can-lose''
        flags when a chained/parity rule depends on deliberately losing?
\end{itemize}
\end{stuckbox}

\begin{practiceset}
  \practice{Serval and Rooted Tree (minimax as a sacrifice count)}{Codeforces 1153D}
    {https://codeforces.com/problemset/problem/1153/D}{CF 2000 · L5}
    {A rooted tree where internal nodes are labelled $\max$ or $\min$; assign $1 \dots k$
    to the $k$ leaves (each once) to \emph{maximise} the root. It is a two-player-style
    $\max/\min$ fold, but the neat resolution is a counting dual: let $d[v] = $ how many
    leaves under $v$ must be ``given big numbers''; a $\min$ node \emph{sums} its children
    (all must be high) and a $\max$ node takes the $\min$ over children (one high branch
    suffices), leaf $d = 1$; the answer is $k + 1 - d[\text{root}]$. Trace the four
    official samples $1, 4, 4, 5$ before trusting the operator -- the swap (min-sums,
    max-mins) is the whole problem.}\end{practiceset}

\paragraph{Back-cites (URLs live elsewhere in the book).}
\begin{itemize}[nosep,leftmargin=1.4em]
  \item \textbf{CSES 2103 / 2104 (Nim, Stick Games)} -- classic $\operatorname{mex}$ grounds for
        \thmref{c27-grundy-xor} on non-tree boards; if you have not seen $\operatorname{mex}$ outside a
        tree, build the nim-heap intuition there first (their URLs are unclaimed in this
        book, so consult CSES directly).
  \item \textbf{Ch.~13 \secref{c10-twostate}}: the in/out Boolean fold is the exact
        ancestor of the $\mathit{win}/\mathit{lose}$ two-flag DP here -- same shape, the
        $\bigvee\neg$ recurrence instead of $\max$. Read it for the base-case discipline.
  \item \textbf{Ch.~15 (\chref{c20})}: rerooting the cutting game over every possible
        ground node is a forward obligation; only the per-root fold
        \coderef{c27-grundy} is used here.
\end{itemize}

\begin{recallbox}[10-second refresher -- 14.4 Games on trees]
A tree game is a bottom-up fold with a game operator. General impartial game:
$g(v) = \operatorname{mex}$ of the option values, leaf $g=0$; disjoint branches combine by
\textbf{XOR} (\thmref{c27-grundy-xor}); cutting games give the closed form
$g(v) = \bigoplus_c (g(c)+1)$ (\thmref{c27-hackenbush}). First wins iff the root value is
$\ne 0$. On $T$, $g(1) = 4 \oplus 3 = 7$, so First wins. When you only need win/lose, drop
to two flags $\mathit{win}[v] = \bigvee_c \neg\mathit{win}[c]$, $\mathit{lose}[v] =
\bigvee_c \neg\mathit{lose}[c]$ -- CF~455B (\probref{c27-game-solved}): verdict $= \neg W
\Rightarrow$ Second, $W \wedge L \Rightarrow$ First, else parity of $k$; samples First,
First, Second.
\selftest{
  \item Why does a two-branch node with each branch value~$2$ become a losing position?
  \item State $\thmref{c27-hackenbush}$'s $+1$: which edge does the $1$ represent?
  \item For CF~455B, give the leaf base $(\mathit{win}, \mathit{lose})$ and why it is
        $(0, 1)$ not $(0, 0)$.
  \item When do you use two Boolean flags instead of a full Grundy number?
  \item On $T$, compute $g(3)$ and $g(1)$ from scratch and read off the winner.
}
\end{recallbox}
